\documentclass[a4paper,11pt,reqno]{amsart}

% =========================================================================
%  Instrumentación Electrónica — PreLab de la sesión 2 (filtro activo de segundo orden)
%  Mismo formato que el PreLab de la sesión 1, con la paleta crema/lila/índigo/noche.
%  Compilar con: latexmk -pdf main.tex. Ecuaciones numeradas de (1) en adelante, a la derecha.
% =========================================================================
% =========================================================================
%  Instrumentación Electrónica — PreLab de la sesión 2 (filtro activo de segundo orden)
%  Estilo heredado de los apuntes de Electromagnetismo I (Temas 0 y 1):
%  amsart, Palatino, versalitas, enunciados numerados por sección, paleta de
%  azules y grises sobrios. Ecuaciones numeradas a la derecha (opción reqno).
% =========================================================================

% --- IDIOMA Y CODIFICACIÓN ------------------------------------------------
\usepackage[utf8]{inputenc}
\usepackage[T1]{fontenc}
\usepackage[spanish,es-tabla,es-nodecimaldot]{babel}

% --- TIPOGRAFÍA (Palatino) --------------------------------------------------
\usepackage{mathpazo}
\linespread{1.05}
\emergencystretch=2em
\usepackage[final]{microtype}

% --- MATEMÁTICAS -------------------------------------------------------------
\usepackage{amsmath, amssymb, mathtools}
\usepackage{bm}

% --- UNIDADES (coma decimal, como en el resto de los apuntes) ---------------
\usepackage{siunitx}
\sisetup{output-decimal-marker={,}, group-separator={\,}, group-minimum-digits=5,
         exponent-product=\cdot, per-mode=symbol, range-phrase={--},
         list-final-separator={ y }, list-pair-separator={ y }}

% --- PÁGINA, COLORES Y GRÁFICOS --------------------------------------------
\usepackage[a4paper,margin=2.6cm]{geometry}
\usepackage{graphicx}
\usepackage{xcolor}
% Paleta de la sesión 2 (crema, lila grisáceo, índigo, azul noche y tinta)
\definecolor{crema}{HTML}{F1DAC4}
\definecolor{lila}{HTML}{A69CAC}
\definecolor{indigo}{HTML}{474973}
\definecolor{noche}{HTML}{161B33}
\definecolor{tinta}{HTML}{0D0C1D}
% Nombres funcionales (los mismos que en la sesión 1, con la nueva paleta)
\definecolor{azulprof}{HTML}{161B33}    % títulos y rótulos (noche)
\definecolor{azulacero}{HTML}{474973}   % enlaces, filetes y curvas (índigo)
\definecolor{azulclaro}{HTML}{A69CAC}   % marcos (lila)
\definecolor{azulfondo}{HTML}{FAF2EA}   % fondo tenue (crema aclarada)
\definecolor{grispizarra}{HTML}{6E6680} % gris lila oscuro (notas)
\definecolor{grisclaro}{HTML}{E4DFE6}   % rejillas
\definecolor{acento}{HTML}{B8936A}      % acento: crema tostada
\definecolor{rojoteo}{HTML}{B8936A}     % resultados y curvas secundarias
\definecolor{verdeteo}{HTML}{474973}    % corrientes
\definecolor{opfondo}{HTML}{F8EDE2}     % relleno del triángulo del AO

\usepackage{tikz}
\usetikzlibrary{babel, arrows.meta, calc, positioning, decorations.markings,
                patterns, intersections}
\usepackage[european, straightvoltages, RPvoltages]{circuitikz}
\usepackage{pgfplots}
\pgfplotsset{compat=1.17}
\usepgfplotslibrary{fillbetween}
\usepackage{booktabs}
\usepackage{needspace}
\usepackage{array}
\usepackage{multirow}
\usepackage[inline]{enumitem}
\setlist[itemize]{leftmargin=1.6em, itemsep=2pt, topsep=4pt}
\setlist[enumerate]{leftmargin=1.8em, itemsep=2pt, topsep=4pt}

% --- ENLACES Y REFERENCIAS CRUZADAS ----------------------------------------
\usepackage[colorlinks=true, linkcolor=azulacero, citecolor=azulacero,
            urlcolor=azulacero, plainpages=false, pdfusetitle=false,
            pdftitle={PreLab. Sesión 2: filtro activo de segundo orden},
            pdfauthor={Daniel Martín Guerrero}]{hyperref}
\usepackage[spanish,nameinlink]{cleveref}

% --- CABECERAS ---------------------------------------------------------------
\pagestyle{headings}
\makeatletter
\def\ps@headings{%
  \let\@mkboth\markboth
  \def\@evenhead{\small\scshape\normalfont\thepage\hfil Instrumentación Electrónica}%
  \def\@oddhead{\small\scshape\normalfont PreLab. Sesión 2: filtro de segundo orden\hfil\thepage}%
  \let\@evenfoot\@empty \let\@oddfoot\@empty}
\makeatother

% Subsecciones como títulos en bloque (no «run-in»), con aire antes y después
\usepackage{titlesec}
\titleformat{\subsection}[block]{\normalfont\bfseries\color{azulprof}}{\thesubsection.}{0.6em}{}
\titlespacing*{\subsection}{0pt}{3ex plus 1ex minus .5ex}{1.6ex plus .4ex}
\setcounter{tocdepth}{2}
\makeatletter
\def\@@and{y}
\renewcommand{\keywordsname}{Palabras clave}
% Índice: número de sección en una caja de anchura fija (evita el solape de amsart)
\renewcommand{\tocsection}[3]{\indentlabel{\@ifnotempty{#2}{\makebox[1.6em][l]{#2.}}}#3}
\renewcommand{\tocsubsection}[3]{\indentlabel{\@ifnotempty{#2}{\makebox[2.4em][l]{#2.}}}#3}
\def\l@section{\@tocline{1}{2pt}{0pt}{1.6em}{}}
\def\l@subsection{\@tocline{2}{0pt}{1.6em}{2.4em}{}}
\def\l@part{\@tocline{-1}{8pt plus 2pt}{0pt}{0pt}{\bfseries\color{azulprof}}}
\makeatother

% =========================================================================
%               NOTACIÓN Y MACROS DE USO FRECUENTE
% =========================================================================
\newcommand{\dif}{\mathrm{d}}
\newcommand{\eul}{\mathrm{e}}
\newcommand{\jj}{\mathrm{j}}                         % unidad imaginaria (j)
\newcommand{\R}{\mathbb{R}}
\newcommand{\C}{\mathbb{C}}
\newcommand{\abs}[1]{\lvert #1\rvert}
\newcommand{\defeq}{\mathrel{\overset{\mathrm{def}}{=}}}
\newcommand{\Lt}{\mathcal{L}}                        % transformada de Laplace
\DeclareMathOperator{\sen}{sen}
\DeclareMathOperator{\sgn}{sgn}
\DeclareMathOperator{\arctg}{arctg}
\DeclareMathOperator{\Real}{Re}
\renewcommand{\sin}{\sen}                            % «sen» en todo el texto
% Magnitudes de circuito
\newcommand{\vin}{v_{\mathrm{in}}}
\newcommand{\vout}{v_{\mathrm{out}}}
\newcommand{\Vin}{V_{\mathrm{in}}}
\newcommand{\Vout}{V_{\mathrm{out}}}
\newcommand{\Vsat}{V_{\mathrm{sat}}}
\newcommand{\Vcc}{V_{\mathrm{CC}}}
\newcommand{\Vee}{V_{\mathrm{EE}}}
\newcommand{\Rin}{R_{\mathrm{in}}}
\newcommand{\Rout}{R_{\mathrm{out}}}
\newcommand{\Zin}{Z_{\mathrm{in}}}
\newcommand{\Zout}{Z_{\mathrm{out}}}
\newcommand{\Av}{A_{V}}
\newcommand{\vp}{v_{+}}
\newcommand{\vm}{v_{-}}
\newcommand{\dB}{\,\mathrm{dB}}
\newcommand{\pend}[1]{\textcolor{grispizarra}{\colorbox{grisclaro}{\strut\,#1\,}}}

% =========================================================================
%        ENTORNOS FORMALES (enunciados numerados, estilo artículo)
% =========================================================================
\usepackage{amsthm}
\usepackage[most]{tcolorbox}

\newtheoremstyle{cienplano}{10pt}{10pt}{\itshape}{}%
  {\bfseries\color{azulprof}}{.}{0.6em}%
  {\thmname{#1}\thmnumber{ #2}\thmnote{\normalfont\upshape\ (#3)}}
\newtheoremstyle{ciendef}{10pt}{10pt}{\normalfont}{}%
  {\bfseries\color{azulprof}}{.}{0.6em}%
  {\thmname{#1}\thmnumber{ #2}\thmnote{\normalfont\upshape\ (#3)}}
\newtheoremstyle{cienobs}{8pt}{8pt}{\normalfont}{}%
  {\itshape\color{azulprof}}{.}{0.6em}%
  {\thmname{#1}\thmnumber{ #2}\thmnote{\normalfont\upshape\ (#3)}}

\theoremstyle{cienplano}
\newtheorem{proposicion}{Proposición}
\newtheorem{teorema}[proposicion]{Teorema}
\newtheorem{corolario}[proposicion]{Corolario}
\theoremstyle{ciendef}
\newtheorem{ejemplo}[proposicion]{Ejemplo}
\newtheorem{ejercicio}[proposicion]{PreLab}
\newtheorem{hipotesis}[proposicion]{Hipótesis}
% Definiciones: único entorno enmarcado (caja fina, fondo tenue, título en el borde)
\newcounter{definicion}
\renewcommand{\thedefinicion}{\theproposicion}
\newtcolorbox{definicionbox}[1]{%
  enhanced, breakable, sharp corners,
  colback=azulfondo, colframe=azulclaro, boxrule=0.4pt,
  left=9pt, right=9pt, top=5pt, bottom=5pt,
  before skip=12pt, after skip=12pt,
  attach boxed title to top left={yshift*=-\tcboxedtitleheight/2, xshift=8pt},
  boxed title style={colback=white, colframe=white, sharp corners, boxrule=0pt,
                     left=3pt, right=3pt, top=1pt, bottom=1pt},
  fonttitle=\bfseries\color{azulprof}, coltitle=azulprof,
  title={Definición~\thedefinicion\ifx\relax#1\relax\else\ (#1)\fi}}
\newenvironment{definicion}[1][]
  {\par\stepcounter{proposicion}\refstepcounter{definicion}\begin{definicionbox}{#1}}
  {\end{definicionbox}}

% Enunciado de cada recuadro del PreLab: misma caja, con filete ámbar (los
% recuadros amarillos del guion), para distinguir lo pedido de lo resuelto.
\newcounter{prelab}
\newtcolorbox{prelabbox}[1]{%
  enhanced, sharp corners,
  colback=white, colframe=crema!90!acento, boxrule=0.5pt, leftrule=2.2pt,
  left=9pt, right=9pt, top=5pt, bottom=5pt,
  before skip=12pt, after skip=10pt,
  attach boxed title to top left={yshift*=-\tcboxedtitleheight/2, xshift=8pt},
  boxed title style={colback=white, colframe=white, sharp corners, boxrule=0pt,
                     left=3pt, right=3pt, top=1pt, bottom=1pt},
  fonttitle=\bfseries\color{indigo}, coltitle=indigo,
  title={PreLab~\theprelab\ifx\relax#1\relax\else\ (#1)\fi}}
\newenvironment{prelab}[1][]
  {\par\refstepcounter{prelab}\begin{prelabbox}{#1}}
  {\end{prelabbox}}

\theoremstyle{cienobs}
\newtheorem{observacion}[proposicion]{Observación}

% Resolución: misma maquinaria que «proof», con otro rótulo
\renewcommand{\qedsymbol}{\ensuremath{\blacksquare}}
\newenvironment{resolucion}[1][Resolución]
  {\begin{proof}[\itshape\color{azulprof}#1]}{\end{proof}}

% Resultado final enmarcado (como \boxed, pero con el filete de la paleta)
\newcommand{\resultado}[1]{\fcolorbox{azulacero}{white}{\ensuremath{\displaystyle #1}}}

\crefname{proposicion}{proposición}{proposiciones}
\Crefname{proposicion}{Proposición}{Proposiciones}
\crefname{teorema}{teorema}{teoremas}
\Crefname{teorema}{Teorema}{Teoremas}
\crefname{corolario}{corolario}{corolarios}
\Crefname{corolario}{Corolario}{Corolarios}
\crefname{definicion}{definición}{definiciones}
\Crefname{definicion}{Definición}{Definiciones}
\crefname{ejemplo}{ejemplo}{ejemplos}
\Crefname{ejemplo}{Ejemplo}{Ejemplos}
\crefname{hipotesis}{hipótesis}{hipótesis}
\Crefname{hipotesis}{Hipótesis}{Hipótesis}
\crefname{prelab}{PreLab}{PreLabs}
\Crefname{prelab}{PreLab}{PreLabs}
\crefname{observacion}{observación}{observaciones}
\Crefname{observacion}{Observación}{Observaciones}
\crefname{section}{sección}{secciones}
\Crefname{section}{Sección}{Secciones}
\crefname{subsection}{sección}{secciones}
\Crefname{subsection}{Sección}{Secciones}
\crefname{table}{tabla}{tablas}
\Crefname{table}{Tabla}{Tablas}
\crefname{figure}{figura}{figuras}
\Crefname{figure}{Figura}{Figuras}
\crefname{equation}{ecuación}{ecuaciones}
\Crefname{equation}{Ecuación}{Ecuaciones}
\crefformat{equation}{ecuación~(#2#1#3)}
\Crefformat{equation}{Ecuación~(#2#1#3)}
\crefrangeformat{equation}{ecuaciones~(#3#1#4)--(#5#2#6)}
\crefmultiformat{equation}{ecuaciones~(#2#1#3)}{ y~(#2#1#3)}{, (#2#1#3)}{ y~(#2#1#3)}

\usepackage[font=small, labelfont={sc,color=azulprof}, labelsep=period,
            margin=0.6cm, justification=justified]{caption}
\usepackage{subcaption}

% Referencia a un panel de figura: «figura 5\,(b)», enlazada en bloque
\newcommand{\figpanel}[2]{\hyperref[#1]{figura~\ref*{#1}\,\textup{(#2)}}}
\newcommand{\Figpanel}[2]{\hyperref[#1]{Figura~\ref*{#1}\,\textup{(#2)}}}

% =========================================================================
%        ESTILOS GRÁFICOS COMUNES (TikZ / circuitikz / pgfplots)
% =========================================================================
% =========================================================================
%  Estilos gráficos comunes: TikZ, circuitikz y pgfplots
% =========================================================================

% --- Estilos TikZ generales ------------------------------------------------
\tikzset{
  >=Stealth,
  eje/.style={->, line width=0.6pt, grispizarra},
  fino/.style={line width=0.4pt, grispizarra},
  punteado/.style={densely dotted, line width=0.5pt, grispizarra},
  etiq/.style={font=\small, text=black!85},
  nota/.style={font=\scriptsize, text=grispizarra, align=center},
  % corriente dibujada junto a un hilo (flecha verde con rótulo)
  corr/.style={-{Stealth[length=4.5pt,width=4pt]}, line width=0.8pt, verdeteo},
  % nudo con tensión anotada
  tens/.style={font=\small, text=azulprof},
}

% --- circuitikz: trazo y colores sobrios -----------------------------------
\ctikzset{
  bipoles/thickness=1,
  bipoles/length=1.15cm,
  tripoles/op amp/width=1.5,
  tripoles/op amp/height=1.35,
  tripoles/op amp/font=\small,
  current/distance=0.45,
  label/align=smart,
}
\tikzset{
  circuito/.style={line width=0.55pt, color=black!85, font=\small},
  ao/.style={fill=opfondo},
}

% --- Estilo base de ejes pgfplots ------------------------------------------
\pgfplotsset{
  tema/.style={
    tick label style={font=\scriptsize}, label style={font=\small},
    axis lines=left, axis line style={grispizarra},
    grid=major, grid style={grisclaro, line width=0.4pt},
    every tick/.style={grispizarra},
    legend style={font=\scriptsize, draw=none, fill=white, fill opacity=0.85,
                  text opacity=1},
    legend cell align=left,
    /pgf/number format/use comma,
  },
  % ejes centrados (formas de onda y características de transferencia)
  temacentro/.style={
    tick label style={font=\scriptsize}, label style={font=\small},
    axis lines=middle, axis line style={grispizarra, -{Stealth[length=4pt]}},
    every tick/.style={grispizarra},
    legend style={font=\scriptsize, draw=none, fill=white, fill opacity=0.85,
                  text opacity=1},
    legend cell align=left,
    /pgf/number format/use comma,
  },
  senal entrada/.style={azulacero, line width=1.0pt},
  senal salida/.style={rojoteo, line width=1.1pt},
  asintota/.style={grispizarra, densely dashed, line width=0.6pt},
}


% Texto principal en tinta
\AtBeginDocument{\color{tinta}}


\newcommand{\Universidad}{Universidad de Salamanca}
\newcommand{\Facultad}{Facultad de Ciencias}
\newcommand{\Titulacion}{Grado en Física}
\newcommand{\Asignatura}{Instrumentación Electrónica}
\newcommand{\NumTema}{Laboratorio de circuitos, sesión 2}
\newcommand{\TituloTema}{PreLab:\\filtro de segundo orden}
\newcommand{\tipotrabajo}{PreLab. Cálculos de diseño}
\newcommand{\miNombre}{Daniel Martín Guerrero}
\newcommand{\miRotuloProfesor}{Profesores}
\newcommand{\miProfesor}{Yahya Moubarak Meziani\\María Susana Pérez Santos\\Jesús Enrique Velázquez Pérez}
\newcommand{\miDetalleProfesor}{Área de Electrónica}
\newcommand{\SubtituloTema}{Función de transferencia, parámetros $T_0$, $\omega_0$ y $Q$\\y respuesta en frecuencia de un filtro activo paso bajo}
\newcommand{\miUbicacion}{Salamanca, curso 2026/2027}

\title{PreLab de la sesión 2: filtro activo paso bajo de segundo orden}
\author{Daniel Martín Guerrero}
\address{Grado en Física, Facultad de Ciencias, Universidad de Salamanca,
         Salamanca, España}
\date{}

\begin{document}
\pagenumbering{alph}
%%%%%%%%%%%%%%%%%%%%%%%%%%%%%%%%%%%%%%%%%%%%%%%%%%%%%%%%%%%%%%%%%%%%%%%%
% Portada formal (estilo tesis): logotipos, identificación institucional,
% título y datos de autoría. Banda vertical a la derecha, como en el
% «Cálculo integral multivariable» de L. Navas. Requiere en main.tex las
% macros \Universidad, \Facultad, \Titulacion, \Asignatura, \NumTema,
% \TituloTema, \SubtituloTema, \tipotrabajo, \miNombre, \miProfesor,
% \miUbicacion, \miRotuloProfesor, \miDetalleProfesor.
%%%%%%%%%%%%%%%%%%%%%%%%%%%%%%%%%%%%%%%%%%%%%%%%%%%%%%%%%%%%%%%%%%%%%%%%

\begin{titlepage}
\thispagestyle{empty}
\pagecolor{white}

\begin{tikzpicture}[remember picture, overlay]
  \coordinate (O) at (current page.south west);

  % Velo muy tenue en la mitad superior (cabecera institucional)
  \shade[top color=crema!45, bottom color=white]
    ($(O)+(0,18.5cm)$) rectangle ($(O)+(21cm,29.7cm)$);

  % Banda vertical azul, a la derecha del bloque de texto
  \fill[indigo] ($(O)+(19.0cm,2.6cm)$) rectangle ($(O)+(19.14cm,27.4cm)$);

  % Logotipos
  \node[anchor=north west, inner sep=0] at ($(O)+(2.4cm,27.7cm)$)
    {\includegraphics[height=1.55cm]{logo_usal.png}};
  \node[anchor=north east, inner sep=0] at ($(O)+(17.4cm,27.7cm)$)
    {\includegraphics[height=1.55cm]{logo_facultad.png}};
  \draw[grispizarra!60, line width=0.4pt] ($(O)+(2.4cm,25.7cm)$) -- ($(O)+(17.4cm,25.7cm)$);

  % Identificación institucional
  \node[anchor=north, align=center, inner sep=0, text width=15cm] at ($(O)+(9.9cm,25.2cm)$)
    {{\color{azulprof}\scshape\Large\textls[120]{\Universidad}}\\[5pt]
     {\color{grispizarra}\scshape\large\textls[80]{\Facultad}}\\[3pt]
     {\color{grispizarra}\scshape\large\textls[80]{\Titulacion}}};

  % Asignatura y título
  \node[anchor=north, align=center, inner sep=0, text width=15cm] at ($(O)+(9.9cm,19.2cm)$)
    {{\color{azulacero}\scshape\Large\textls[160]{\Asignatura}}\\[4pt]
     {\color{grispizarra}\scshape\large\textls[100]{\NumTema}}\\[22pt]
     {\color{azulprof}\fontsize{30}{36}\selectfont\bfseries\TituloTema}\\[14pt]
     {\color{azulacero}\rule{3.2cm}{0.8pt}}\\[14pt]
     {\color{tinta!85}\itshape\large\SubtituloTema}};

  % Epígrafe: forma canónica del filtro paso bajo de segundo orden
  \node[anchor=north, align=center, inner sep=0, text width=15cm] at ($(O)+(9.9cm,11.0cm)$)
    {\color{grispizarra}\small
     $\displaystyle T(s)=\frac{T_{0}}{s^{2}+\dfrac{\omega_{0}}{Q}\,s+\omega_{0}^{2}}\qquad
      \omega_{0}=\frac{1}{\sqrt{R_{1}R_{2}C_{1}C_{2}}}\qquad
      K=1+\frac{R_{A}}{R_{B}}$};

  % Autoría y dirección
  \draw[grispizarra!60, line width=0.4pt] ($(O)+(2.4cm,8.6cm)$) -- ($(O)+(17.4cm,8.6cm)$);
  \node[anchor=north west, inner sep=0, text width=7cm, align=left] at ($(O)+(2.4cm,8.2cm)$)
    {{\scshape\small\color{azulacero}\textls[80]{Autor}}\\[4pt]
     {\large\color{tinta}\miNombre}\\[2pt]
     {\small\color{grispizarra}Grado en Física}};
  \node[anchor=north west, inner sep=0, text width=7cm, align=left] at ($(O)+(10.4cm,8.2cm)$)
    {{\scshape\small\color{azulacero}\textls[80]{\miRotuloProfesor}}\\[4pt]
     {\color{tinta}\miProfesor}\\[3pt]
     {\small\color{grispizarra}\miDetalleProfesor}};

  % Tipo de documento, lugar y curso
  \node[anchor=south, align=center, inner sep=0, text width=15cm] at ($(O)+(9.9cm,2.9cm)$)
    {{\color{azulprof}\scshape\large\textls[100]{\tipotrabajo}}\\[4pt]
     {\color{grispizarra}\scshape\small\textls[80]{\miUbicacion}}};
\end{tikzpicture}

\end{titlepage}

\clearpage
\pagenumbering{arabic}
\setcounter{page}{1}

\maketitle

% =========================================================================
\section{Introducción}\label{sec:introduccion}
% =========================================================================
En la sesión anterior vimos que un condensador en la realimentación convierte un amplificador en un
filtro de primer orden, cuya ganancia cae a $-20\dB$/década. En este PreLab damos un paso más:
analizamos un filtro activo paso bajo de \emph{segundo orden}, con dos condensadores y un único
amplificador operacional, cuya caída es el doble de abrupta. Obtendremos su función de transferencia
$T(s)$, la escribiremos en forma canónica para leer sus parámetros $T_0$, $\omega_0$ y $Q$, y estudiaremos
su respuesta en frecuencia: ganancia en continua, módulo, frecuencia de corte y posibles máximos.
Antes, reunimos brevemente las herramientas que vamos a utilizar.

% -------------------------------------------------------------------------
\subsection{Amplificador operacional ideal}
% -------------------------------------------------------------------------
Tratamos el AO como ideal: resistencia de entrada infinita, resistencia de salida nula y ganancia en lazo
abierto $A$ infinita. Si una red lleva parte de la salida a la entrada inversora (realimentación negativa),
la salida permanece finita y de $\vout=A(\vp-\vm)$ se siguen dos reglas:
\begin{itemize}[leftmargin=3em]
  \item[(R1)] no entra corriente por las entradas, $i_+=i_-=0$;
  \item[(R2)] ambas entradas están al mismo potencial, $\vp=\vm$ (cortocircuito virtual).
\end{itemize}

% -------------------------------------------------------------------------
\subsection{Impedancias y función de transferencia}
% -------------------------------------------------------------------------
Con la transformada de Laplace, $F(s)=\int_0^\infty f(t)\,\eul^{-st}\,\dif t$, y condiciones iniciales nulas,
derivar en el tiempo equivale a multiplicar por $s$: $\Lt\{\dif f/\dif t\}=sF(s)$. Así, cada elemento cumple una
«ley de Ohm» $V(s)=Z(s)\,I(s)$, con
\[
  Z_R=R,\qquad Z_C(s)=\frac{1}{sC},
\]
y las leyes de Kirchhoff, (R1) y (R2), por linealidad, valen para las transformadas: el circuito se resuelve
con álgebra. La \emph{función de transferencia} es el cociente $T(s)=\Vout(s)/\Vin(s)$; en régimen sinusoidal de
pulsación $\omega=2\pi f$ basta tomar $s=\jj\omega$, y entonces $\abs{T(\jj\omega)}$ es el cociente de amplitudes de
salida y entrada y $\arg T(\jj\omega)$, su desfase.

% -------------------------------------------------------------------------
\subsection{Filtro paso bajo de segundo orden}
% -------------------------------------------------------------------------
\begin{definicion}[Forma canónica]
Un filtro paso bajo de segundo orden tiene una función de transferencia de la forma
\[
  T(s)=\frac{T_0}{s^{2}+\dfrac{\omega_0}{Q}\,s+\omega_0^{2}},
\]
donde $\omega_0>0$ es la \emph{frecuencia natural} (o característica), $Q>0$ el \emph{factor de calidad} y $T_0$
una constante que fija la ganancia: en continua, $T(0)=T_0/\omega_0^{2}$.
\end{definicion}

Los tres parámetros tienen un significado claro. A baja frecuencia ($\omega\ll\omega_0$) domina $\omega_0^{2}$ en el
denominador y la ganancia es constante, $T_0/\omega_0^{2}$; a alta frecuencia ($\omega\gg\omega_0$) domina $s^{2}$ y
$\abs{T}\simeq T_0/\omega^{2}$, que cae $40\dB$ por década, el doble que un filtro de primer orden. Entre ambos
regímenes, $Q$ controla la forma del codo: los polos,
\[
  s_{\pm}=-\frac{\omega_0}{2Q}\pm\omega_0\sqrt{\frac{1}{4Q^{2}}-1},
\]
tienen parte real negativa (filtro estable) si $Q>0$, y son complejos conjugados si $Q>1/2$. Cuanto mayor es
$Q$, más se acercan al eje imaginario y más pronunciado es el pico de $\abs{T}$ cerca de $\omega_0$; el caso
$Q=1/\sqrt2$ separa las respuestas con pico de las que no lo tienen (filtro de Butterworth, máximamente
plano). Por último, la \emph{frecuencia de corte} $\omega_{-3\mathrm{dB}}$ es aquella en la que $\abs{T}$ cae a
$1/\sqrt2$ de su valor en continua, es decir, $3\dB$ por debajo.

% =========================================================================
\section{Ejercicio 1: función de transferencia}\label{sec:transferencia}
% =========================================================================
\begin{prelab}[Función de transferencia]
Obtener la función de transferencia $T(s)=\Vout/\Vin$ del circuito de la \cref{fig:filtro} y demostrar que
adopta la forma
\[
  T(s)=\frac{\Vout}{\Vin}=\frac{T_0}{\omega_0^{2}+\dfrac{\omega_0\,s}{Q}+s^{2}} ,
\]
asumiendo que el amplificador operacional es ideal. Obtener las expresiones de $T_0$, $\omega_0$ y $Q$ en función
de $R_1$, $C_1$, $R_2$, $C_2$, $R_A$ y $R_B$.
\end{prelab}

% Figura: filtro activo paso bajo de segundo orden (figura 1 del guion), con nudos y corrientes.
\begin{figure}[ht]
\centering
\begin{circuitikz}[circuito]
  \draw (0,0) node[op amp, ao, noinv input up] (A) {};
  % salida
  \draw (A.out) -- ++(0.9,0) coordinate (O);
  \draw (O) to[short, *-o] ++(0.8,0) node[right, tens] {$\Vout$};
  % entrada no inversora: nudo P, C2 a tierra
  \draw (A.+) -- ++(-1.5,0) coordinate (P);
  \draw (P) to[C, l_=$C_2$, i>_=$I_{C2}$] ++(0,-2.4) node[ground] {};
  % R2 entre X y P
  \draw (P) to[R, l_=$R_2$, i<_=$I_2$] ++(-2.2,0) coordinate (X);
  % R1 entre la entrada y X
  \draw (X) to[R, l_=$R_1$, i<^=$I_1$] ++(-2.3,0) coordinate (E);
  \draw (E) to[short, -o] ++(-0.2,0) coordinate (Ein) node[above, tens] {$\Vin$};
  \draw (Ein) -- ++(-0.7,0) to[sV, l_=$v_{\mathrm{in}}$] ++(0,-2.6) node[ground] {};
  % C1 entre X y la salida
  \draw (X) -- ++(0,1.9) coordinate (T) to[C, l=$C_1$, i>^=$I_{C1}$] (T -| O) -- (O);
  % red RA-RB en la entrada inversora
  \draw (A.-) -- ++(-0.4,0) -- ++(0,-1.3) coordinate (N);
  \draw (N) to[R, l_=$R_A$] (N -| O) -- (O);
  \draw (N) to[R, l=$R_B$] ++(0,-1.6) node[ground] {};
  % nudos y tensiones
  \fill (X) circle (1.6pt); \fill (P) circle (1.6pt); \fill (N) circle (1.6pt);
  \node[tens, anchor=north west, font=\scriptsize] at ($(X)+(0.06,-0.06)$) {$V_x$};
  \node[tens, anchor=south west, font=\scriptsize] at ($(P)+(0.06,0.06)$) {$V_+$};
  \node[tens, anchor=north east, font=\scriptsize] at ($(N)+(-0.06,-0.06)$) {$V_-$};
\end{circuitikz}
\caption{Filtro activo paso bajo de segundo orden (figura~1 del guion), con la alimentación omitida.
Se indican los nudos del análisis ($V_x$, $V_+$, $V_-$) y los sentidos de corriente elegidos.}
\label{fig:filtro}
\end{figure}


\begin{resolucion}
Trabajamos en el dominio de Laplace con condiciones iniciales nulas, de modo que los condensadores tienen
impedancias $1/(sC_1)$ y $1/(sC_2)$, y aplicamos al AO ideal las reglas (R1) y (R2); la realimentación negativa
la proporciona la red $R_A$--$R_B$. Planteamos la ley de Kirchhoff de corrientes en los nudos $V_-$, $V_+$ y
$V_x$, con los sentidos de la \cref{fig:filtro}.

\textbf{1. Nudo inversor: divisor $R_A$--$R_B$.} Como $I_-=0$ (R1), por $R_A$ y $R_B$ circula la misma
corriente, y forman un divisor de tensión de la salida:
\begin{equation}\label{eq:vm}
  V_-(s)=\frac{R_B}{R_A+R_B}\,\Vout(s) .
\end{equation}
Por (R2), $V_+=V_-$. Definiendo la ganancia del amplificador no inversor,
\begin{equation}\label{eq:K}
  K\equiv1+\frac{R_A}{R_B}=\frac{R_A+R_B}{R_B},
\end{equation}
la tensión en la entrada no inversora queda ligada a la salida:
\begin{equation}\label{eq:vp}
  V_+(s)=V_-(s)=\frac{\Vout(s)}{K} .
\end{equation}

\textbf{2. Nudo $V_+$.} Como $I_+=0$ (R1), toda la corriente que llega por $R_2$ se va por $C_2$ a tierra:
\begin{equation}\label{eq:nudoP}
  I_2=I_{C2}\quad\Longrightarrow\quad\frac{V_x-V_+}{R_2}=sC_2\,V_+ .
\end{equation}
Despejamos $V_x$ y sustituimos \eqref{eq:vp}:
\begin{equation}\label{eq:vx}
  V_x=(1+sR_2C_2)\,V_+=\frac{1+sR_2C_2}{K}\,\Vout .
\end{equation}

\textbf{3. Nudo $V_x$.} La corriente que entra por $R_1$ sale por $R_2$ y por $C_1$ (hacia la salida):
\begin{equation}\label{eq:nudoX}
  I_1=I_2+I_{C1}\quad\Longrightarrow\quad\frac{\Vin-V_x}{R_1}=\frac{V_x-V_+}{R_2}+sC_1\,(V_x-\Vout) .
\end{equation}
Por \eqref{eq:nudoP} y \eqref{eq:vp}, el primer sumando del segundo miembro vale $sC_2V_+=sC_2\Vout/K$.
Multiplicando por $R_1$ y despejando $\Vin$:
\begin{equation}\label{eq:vin1}
  \Vin=(1+sR_1C_1)\,V_x+\frac{sR_1C_2}{K}\,\Vout-sR_1C_1\,\Vout .
\end{equation}

\textbf{4. Sustitución.} Llevamos \eqref{eq:vx} a \eqref{eq:vin1} y sacamos factor común $\Vout/K$:
\begin{equation}\label{eq:vin2}
  \Vin=\frac{\Vout}{K}\Bigl[(1+sR_1C_1)(1+sR_2C_2)+sR_1C_2-sK R_1C_1\Bigr] .
\end{equation}
Desarrollando el producto, $(1+sR_1C_1)(1+sR_2C_2)=1+s(R_1C_1+R_2C_2)+s^{2}R_1R_2C_1C_2$, y agrupando por
potencias de $s$:
\begin{equation}\label{eq:vin3}
  \Vin=\frac{\Vout}{K}\Bigl[s^{2}R_1R_2C_1C_2+s\bigl(R_1C_2+R_2C_2+(1-K)R_1C_1\bigr)+1\Bigr] .
\end{equation}

\textbf{5. Función de transferencia.} Dividiendo $\Vout$ entre $\Vin$ en \eqref{eq:vin3}:
\begin{equation}\label{eq:T1}
  T(s)=\frac{\Vout}{\Vin}=\frac{K}{s^{2}R_1R_2C_1C_2+s\bigl(R_1C_2+R_2C_2+(1-K)R_1C_1\bigr)+1} .
\end{equation}
Para llevarla a la forma pedida, dividimos numerador y denominador por $R_1R_2C_1C_2$, de modo que el
coeficiente de $s^{2}$ sea $1$:
\begin{equation}\label{eq:T2}
  \resultado{T(s)=\frac{\dfrac{K}{R_1R_2C_1C_2}}{s^{2}+s\Bigl(\dfrac{1}{R_1C_1}+\dfrac{1}{R_2C_1}+\dfrac{1-K}{R_2C_2}\Bigr)+\dfrac{1}{R_1R_2C_1C_2}}}
\end{equation}
que tiene exactamente la forma del enunciado: una constante en el numerador y un polinomio de segundo grado
en $s$ en el denominador. Es, por tanto, un filtro paso bajo de segundo orden.

\textbf{6. Identificación de los parámetros.} Comparando \eqref{eq:T2} término a término con la forma
del enunciado, $T_0/(s^{2}+\omega_0 s/Q+\omega_0^{2})$:
\begin{align}
  \omega_0^{2}&=\frac{1}{R_1R_2C_1C_2}
  &&\Longrightarrow\quad \resultado{\omega_0=\frac{1}{\sqrt{R_1R_2C_1C_2}}}\label{eq:w0}\\
  T_0&=\frac{K}{R_1R_2C_1C_2}=K\,\omega_0^{2}
  &&\Longrightarrow\quad \resultado{T_0=\frac{1+R_A/R_B}{R_1R_2C_1C_2}}\label{eq:T0}\\
  \frac{\omega_0}{Q}&=\frac{1}{R_1C_1}+\frac{1}{R_2C_1}+\frac{1-K}{R_2C_2}. &&\label{eq:w0Q}
\end{align}
Despejando $Q$ de \eqref{eq:w0Q} con \eqref{eq:w0}, y multiplicando numerador y denominador por
$R_1R_2C_1C_2$:
\begin{equation}\label{eq:Q}
  Q=\frac{\omega_0}{\dfrac{1}{R_1C_1}+\dfrac{1}{R_2C_1}+\dfrac{1-K}{R_2C_2}}
   =\frac{\sqrt{R_1R_2C_1C_2}}{R_1C_2+R_2C_2+(1-K)\,R_1C_1} ,
\end{equation}
y, como $1-K=-R_A/R_B$ por \eqref{eq:K},
\begin{equation}\label{eq:Qfinal}
  \resultado{Q=\frac{\sqrt{R_1R_2C_1C_2}}{(R_1+R_2)\,C_2-\dfrac{R_A}{R_B}\,R_1C_1}}
\end{equation}

\textbf{7. Comprobaciones y comentarios.}
\begin{itemize}
  \item \emph{Límites físicos.} Para $s\to0$ los condensadores son circuitos abiertos: no circula corriente por
        $R_1$ ni $R_2$, $V_+=\Vin$ y el circuito es un no inversor, $\Vout=K\Vin$. En efecto, \eqref{eq:T1} da
        $T(0)=K$. Para $s\to\infty$, $C_2$ cortocircuita $V_+$ a tierra y $\Vout\to0$; en efecto,
        $T(s)\simeq T_0/s^{2}\to0$.
  \item \emph{Papel de cada componente.} $\omega_0$ \eqref{eq:w0} solo depende de las resistencias y
        condensadores de la red RC; la ganancia $K$ solo entra en $T_0$ \eqref{eq:T0} y en $Q$ \eqref{eq:Qfinal}.
        Aumentar $R_A/R_B$ reduce el denominador de \eqref{eq:Qfinal} y, por tanto, \emph{aumenta} $Q$: $C_1$
        devuelve parte de la salida al nudo $V_x$ (realimentación positiva), que realza la respuesta cerca de
        $\omega_0$.
  \item \emph{Estabilidad.} Los polos tienen parte real negativa si $Q>0$, es decir, si el denominador de
        \eqref{eq:Qfinal} es positivo: $R_A/R_B<(R_1+R_2)C_2/(R_1C_1)$. Si se supera ese valor, el circuito
        deja de ser un filtro y oscila.
\end{itemize}
\end{resolucion}

% =========================================================================
\section{Ejercicio 2: respuesta en frecuencia}\label{sec:respuesta}
% =========================================================================
\begin{prelab}[Respuesta en frecuencia]
Obténgase: (a) el valor de la función de transferencia para $\omega=0$; (b) la expresión del módulo de
$T(s)$; (c) la frecuencia $\omega_{-3\mathrm{dB}}$. ¿Cuántos máximos tiene la función $\abs{T(s)}$? ¿A qué
frecuencia(s) está(n) localizado(s) y qué valor(es) adopta $\abs{T(s)}$ a esas frecuencias?
\end{prelab}

\begin{resolucion}
Partimos de la forma canónica obtenida en el ejercicio 1, \eqref{eq:T2}, con los parámetros \eqref{eq:w0},
\eqref{eq:T0} y \eqref{eq:Qfinal}. En régimen sinusoidal tomamos $s=\jj\omega$; como $s^{2}=-\omega^{2}$,
\begin{equation}\label{eq:Tjw}
  T(\jj\omega)=\frac{T_0}{(\omega_0^{2}-\omega^{2})+\jj\,\dfrac{\omega_0\,\omega}{Q}} .
\end{equation}

\textbf{(a) Ganancia en continua.} Con $\omega=0$ en \eqref{eq:Tjw} y $T_0=K\omega_0^{2}$ \eqref{eq:T0}:
\begin{equation}\label{eq:T0dc}
  T(0)=\frac{T_0}{\omega_0^{2}}\quad\Longrightarrow\quad\resultado{T(0)=K=1+\frac{R_A}{R_B}}
\end{equation}
Es real y positiva: en continua los condensadores son circuitos abiertos y el circuito se reduce a un
amplificador no inversor, sin desfase.

\textbf{(b) Módulo.} El módulo de un cociente es el cociente de los módulos, y $\abs{a+\jj b}=\sqrt{a^{2}+b^{2}}$:
\begin{equation}\label{eq:mod}
  \resultado{\abs{T(\jj\omega)}=\frac{T_0}{\sqrt{(\omega_0^{2}-\omega^{2})^{2}+\Bigl(\dfrac{\omega_0\,\omega}{Q}\Bigr)^{2}}}}
\end{equation}
Conviene normalizarlo. Dividiendo numerador y denominador por $\omega_0^{2}$, con $T_0/\omega_0^{2}=K$ y la
pulsación reducida $u=\omega/\omega_0$:
\begin{equation}\label{eq:modn}
  \abs{T(\jj\omega)}=\frac{K}{\sqrt{(1-u^{2})^{2}+u^{2}/Q^{2}}},\qquad u=\frac{\omega}{\omega_0} .
\end{equation}
Para $u\ll1$, $\abs{T}\to K$; para $u\gg1$, $\abs{T}\simeq K/u^{2}$, que cae $40\dB$ por década: es un paso bajo de
segundo orden. En $u=1$, $\abs{T(\jj\omega_0)}=KQ$.

\textbf{(c) Frecuencia de corte.} La definimos como la pulsación a la que el módulo cae a $1/\sqrt2$ de su
valor en continua \eqref{eq:T0dc}. Igualando \eqref{eq:modn} a $K/\sqrt2$ y elevando al cuadrado:
\begin{equation}\label{eq:c1}
  (1-u^{2})^{2}+\frac{u^{2}}{Q^{2}}=2 .
\end{equation}
Con $x=u^{2}\ge0$, \eqref{eq:c1} es una ecuación de segundo grado:
\begin{equation}\label{eq:c2}
  x^{2}-\Bigl(2-\frac{1}{Q^{2}}\Bigr)x-1=0 .
\end{equation}
El producto de sus raíces es $-1$, de modo que tiene una raíz positiva y otra negativa; solo la positiva es
físicamente válida:
\begin{equation}\label{eq:c3}
  x_{-3\mathrm{dB}}=1-\frac{1}{2Q^{2}}+\sqrt{\Bigl(1-\frac{1}{2Q^{2}}\Bigr)^{2}+1} .
\end{equation}
Deshaciendo el cambio, $\omega=\omega_0\sqrt{x}$:
\begin{equation}\label{eq:w3dB}
  \resultado{\omega_{-3\mathrm{dB}}=\omega_0\,\sqrt{1-\frac{1}{2Q^{2}}+\sqrt{\Bigl(1-\frac{1}{2Q^{2}}\Bigr)^{2}+1}}}
\end{equation}
En el caso $Q=1/\sqrt2$, $1-1/(2Q^{2})=0$ y \eqref{eq:w3dB} da $\omega_{-3\mathrm{dB}}=\omega_0$: la frecuencia
natural coincide con la de corte.

\textbf{Máximos de $\abs{T}$.} Por \eqref{eq:modn}, $\abs{T}$ es máximo donde es mínima la función
\begin{equation}\label{eq:D}
  D(x)=(1-x)^{2}+\frac{x}{Q^{2}},\qquad x=u^{2}\ge0 .
\end{equation}
Derivando e igualando a cero:
\begin{equation}\label{eq:dD}
  \frac{\dif D}{\dif x}=-2(1-x)+\frac{1}{Q^{2}}=0\quad\Longrightarrow\quad x_p=1-\frac{1}{2Q^{2}},
\end{equation}
y, como $\dif^{2}D/\dif x^{2}=2>0$, $D$ es una parábola con un único mínimo en $x_p$. Hay dos casos:
\begin{itemize}
  \item \emph{$Q>1/\sqrt2$}: $x_p>0$, y $\abs{T}$ crece desde $K$ hasta un único máximo (pico de resonancia) en
        \begin{equation}\label{eq:wp}
          \resultado{\omega_p=\omega_0\sqrt{1-\frac{1}{2Q^{2}}}}
        \end{equation}
        Sustituyendo $x_p$ en \eqref{eq:D}, $D(x_p)=\dfrac{1}{4Q^{4}}+\dfrac{1}{Q^{2}}-\dfrac{1}{2Q^{4}}=\dfrac{4Q^{2}-1}{4Q^{4}}$,
        y por \eqref{eq:modn}:
        \begin{equation}\label{eq:Tmax}
          \resultado{\abs{T}_{\max}=\frac{K}{\sqrt{D(x_p)}}=\frac{2KQ^{2}}{\sqrt{4Q^{2}-1}}=\frac{KQ}{\sqrt{1-\dfrac{1}{4Q^{2}}}}}
        \end{equation}
        que es mayor que $K$. En $\omega=0$ no hay máximo, porque $D$ todavía decrece ($\dif D/\dif x=-2+1/Q^{2}<0$).
  \item \emph{$Q\le1/\sqrt2$}: $x_p\le0$, de modo que $D$ es creciente para todo $x\ge0$ y $\abs{T}$ decrece
        monótonamente. El único máximo está en $\omega=0$, con valor $\abs{T}_{\max}=T(0)=K$ \eqref{eq:T0dc}.
\end{itemize}

\textbf{Conclusión.} $\abs{T(\jj\omega)}$ tiene \emph{un único máximo}: en $\omega=0$, con valor $K=1+R_A/R_B$, si
$Q\le1/\sqrt2$; o en $\omega_p$ \eqref{eq:wp}, con valor \eqref{eq:Tmax}, si $Q>1/\sqrt2$. El valor
$Q=1/\sqrt2$ es la frontera: la respuesta más plana posible sin pico (\cref{fig:modulo}).
\end{resolucion}

% Figura: módulo normalizado |T|/K frente a omega/omega_0 para varios Q.
\begin{figure}[ht]
\centering
\begin{tikzpicture}
  \begin{semilogxaxis}[tema, width=0.8\textwidth, height=6.2cm,
      xmin=0.1, xmax=10, ymin=-40, ymax=10, ytick={-40,-30,-20,-10,-3,0,6},
      xtick={0.1,1,10}, xticklabels={$0{,}1$,$1$,$10$},
      xlabel={$\omega/\omega_0$}, ylabel={$\abs{T}/K$ (dB)}, log basis x=10,
      legend pos=south west]
    \addplot[noche, line width=1.1pt, domain=0.1:10, samples=250] {-10*log10((1-x^2)^2+x^2/4)};
    \addlegendentry{{$Q=2$}}
    \addplot[indigo, line width=1.1pt, domain=0.1:10, samples=250] {-10*log10((1-x^2)^2+x^2)};
    \addlegendentry{{$Q=1$}}
    \addplot[acento, line width=1.1pt, domain=0.1:10, samples=250] {-10*log10((1-x^2)^2+2*x^2)};
    \addlegendentry{{$Q=1/\sqrt2$}}
    \addplot[lila, line width=1.1pt, domain=0.1:10, samples=250] {-10*log10((1-x^2)^2+4*x^2)};
    \addlegendentry{{$Q=1/2$}}
    \addplot[asintota, domain=0.1:10, samples=2] {-3};
    \addplot[only marks, mark=*, mark size=1.6pt, noche] coordinates {(0.9354,6.30)};
    \addplot[only marks, mark=*, mark size=1.6pt, indigo] coordinates {(0.7071,1.25)};
    \node[font=\scriptsize, anchor=south west, text=grispizarra, rotate=-38] at (axis cs:2.2,-14) {$-40\dB$/década};
  \end{semilogxaxis}
\end{tikzpicture}
\caption{Módulo normalizado $\abs{T}/K$, según \eqref{eq:modn}, para varios valores de $Q$. Para $Q>1/\sqrt2$ aparece
un máximo en $\omega_p<\omega_0$ \eqref{eq:wp}, de altura \eqref{eq:Tmax} (puntos); para $Q\le1/\sqrt2$ el máximo
está en $\omega=0$. Con $Q=1/\sqrt2$ la curva corta la recta de $-3\dB$ justo en $\omega=\omega_0$.}
\label{fig:modulo}
\end{figure}



\end{document}
